% Options for packages loaded elsewhere
\PassOptionsToPackage{unicode}{hyperref}
\PassOptionsToPackage{hyphens}{url}
%
\documentclass[
]{article}
\usepackage{amsmath,amssymb}
\usepackage{iftex}
\ifPDFTeX
  \usepackage[T1]{fontenc}
  \usepackage[utf8]{inputenc}
  \usepackage{textcomp} % provide euro and other symbols
\else % if luatex or xetex
  \usepackage{unicode-math} % this also loads fontspec
  \defaultfontfeatures{Scale=MatchLowercase}
  \defaultfontfeatures[\rmfamily]{Ligatures=TeX,Scale=1}
\fi
\usepackage{lmodern}
\ifPDFTeX\else
  % xetex/luatex font selection
\fi
% Use upquote if available, for straight quotes in verbatim environments
\IfFileExists{upquote.sty}{\usepackage{upquote}}{}
\IfFileExists{microtype.sty}{% use microtype if available
  \usepackage[]{microtype}
  \UseMicrotypeSet[protrusion]{basicmath} % disable protrusion for tt fonts
}{}
\makeatletter
\@ifundefined{KOMAClassName}{% if non-KOMA class
  \IfFileExists{parskip.sty}{%
    \usepackage{parskip}
  }{% else
    \setlength{\parindent}{0pt}
    \setlength{\parskip}{6pt plus 2pt minus 1pt}}
}{% if KOMA class
  \KOMAoptions{parskip=half}}
\makeatother
\usepackage{xcolor}
\usepackage[margin=1in]{geometry}
\usepackage{color}
\usepackage{fancyvrb}
\newcommand{\VerbBar}{|}
\newcommand{\VERB}{\Verb[commandchars=\\\{\}]}
\DefineVerbatimEnvironment{Highlighting}{Verbatim}{commandchars=\\\{\}}
% Add ',fontsize=\small' for more characters per line
\usepackage{framed}
\definecolor{shadecolor}{RGB}{248,248,248}
\newenvironment{Shaded}{\begin{snugshade}}{\end{snugshade}}
\newcommand{\AlertTok}[1]{\textcolor[rgb]{0.94,0.16,0.16}{#1}}
\newcommand{\AnnotationTok}[1]{\textcolor[rgb]{0.56,0.35,0.01}{\textbf{\textit{#1}}}}
\newcommand{\AttributeTok}[1]{\textcolor[rgb]{0.13,0.29,0.53}{#1}}
\newcommand{\BaseNTok}[1]{\textcolor[rgb]{0.00,0.00,0.81}{#1}}
\newcommand{\BuiltInTok}[1]{#1}
\newcommand{\CharTok}[1]{\textcolor[rgb]{0.31,0.60,0.02}{#1}}
\newcommand{\CommentTok}[1]{\textcolor[rgb]{0.56,0.35,0.01}{\textit{#1}}}
\newcommand{\CommentVarTok}[1]{\textcolor[rgb]{0.56,0.35,0.01}{\textbf{\textit{#1}}}}
\newcommand{\ConstantTok}[1]{\textcolor[rgb]{0.56,0.35,0.01}{#1}}
\newcommand{\ControlFlowTok}[1]{\textcolor[rgb]{0.13,0.29,0.53}{\textbf{#1}}}
\newcommand{\DataTypeTok}[1]{\textcolor[rgb]{0.13,0.29,0.53}{#1}}
\newcommand{\DecValTok}[1]{\textcolor[rgb]{0.00,0.00,0.81}{#1}}
\newcommand{\DocumentationTok}[1]{\textcolor[rgb]{0.56,0.35,0.01}{\textbf{\textit{#1}}}}
\newcommand{\ErrorTok}[1]{\textcolor[rgb]{0.64,0.00,0.00}{\textbf{#1}}}
\newcommand{\ExtensionTok}[1]{#1}
\newcommand{\FloatTok}[1]{\textcolor[rgb]{0.00,0.00,0.81}{#1}}
\newcommand{\FunctionTok}[1]{\textcolor[rgb]{0.13,0.29,0.53}{\textbf{#1}}}
\newcommand{\ImportTok}[1]{#1}
\newcommand{\InformationTok}[1]{\textcolor[rgb]{0.56,0.35,0.01}{\textbf{\textit{#1}}}}
\newcommand{\KeywordTok}[1]{\textcolor[rgb]{0.13,0.29,0.53}{\textbf{#1}}}
\newcommand{\NormalTok}[1]{#1}
\newcommand{\OperatorTok}[1]{\textcolor[rgb]{0.81,0.36,0.00}{\textbf{#1}}}
\newcommand{\OtherTok}[1]{\textcolor[rgb]{0.56,0.35,0.01}{#1}}
\newcommand{\PreprocessorTok}[1]{\textcolor[rgb]{0.56,0.35,0.01}{\textit{#1}}}
\newcommand{\RegionMarkerTok}[1]{#1}
\newcommand{\SpecialCharTok}[1]{\textcolor[rgb]{0.81,0.36,0.00}{\textbf{#1}}}
\newcommand{\SpecialStringTok}[1]{\textcolor[rgb]{0.31,0.60,0.02}{#1}}
\newcommand{\StringTok}[1]{\textcolor[rgb]{0.31,0.60,0.02}{#1}}
\newcommand{\VariableTok}[1]{\textcolor[rgb]{0.00,0.00,0.00}{#1}}
\newcommand{\VerbatimStringTok}[1]{\textcolor[rgb]{0.31,0.60,0.02}{#1}}
\newcommand{\WarningTok}[1]{\textcolor[rgb]{0.56,0.35,0.01}{\textbf{\textit{#1}}}}
\usepackage{graphicx}
\makeatletter
\def\maxwidth{\ifdim\Gin@nat@width>\linewidth\linewidth\else\Gin@nat@width\fi}
\def\maxheight{\ifdim\Gin@nat@height>\textheight\textheight\else\Gin@nat@height\fi}
\makeatother
% Scale images if necessary, so that they will not overflow the page
% margins by default, and it is still possible to overwrite the defaults
% using explicit options in \includegraphics[width, height, ...]{}
\setkeys{Gin}{width=\maxwidth,height=\maxheight,keepaspectratio}
% Set default figure placement to htbp
\makeatletter
\def\fps@figure{htbp}
\makeatother
\setlength{\emergencystretch}{3em} % prevent overfull lines
\providecommand{\tightlist}{%
  \setlength{\itemsep}{0pt}\setlength{\parskip}{0pt}}
\setcounter{secnumdepth}{-\maxdimen} % remove section numbering
\ifLuaTeX
  \usepackage{selnolig}  % disable illegal ligatures
\fi
\IfFileExists{bookmark.sty}{\usepackage{bookmark}}{\usepackage{hyperref}}
\IfFileExists{xurl.sty}{\usepackage{xurl}}{} % add URL line breaks if available
\urlstyle{same}
\hypersetup{
  pdftitle={Test},
  pdfauthor={Test},
  hidelinks,
  pdfcreator={LaTeX via pandoc}}

\title{Test}
\author{Test}
\date{8/5/2021}

\begin{document}
\maketitle

\hypertarget{inverse-transform-method-for-sampling-from-random-variable-with-kernel-g}{%
\subsection{\texorpdfstring{Inverse transform method for sampling from
random variable with kernel
\(g\)}{Inverse transform method for sampling from random variable with kernel g}}\label{inverse-transform-method-for-sampling-from-random-variable-with-kernel-g}}

First, we can use the inverse-transform method. The kernel of \(g\) is
\(\exp(-|x|)\), and so the normalizing constant \(Z\) is given by \[
  \frac{1}{Z} \int_{-\infty}^{\infty} g(x) dx = 1,
\] which may be rewritten as \[
  \int_{-\infty}^{0} e^{x} dx + \int_{0}^{\infty} e^{-x} dx = Z,
\] which has the solution \(Z=2\). Thus, random variable \(Y\) that has
a kernel \(g\) has a density \[
  f_Y(y) = \frac{1}{2} \exp(-|x|).
\]

\begin{Shaded}
\begin{Highlighting}[]
\NormalTok{g.density }\OtherTok{\textless{}{-}} \ControlFlowTok{function}\NormalTok{(x) \{ }\FloatTok{0.5}\SpecialCharTok{*}\FunctionTok{exp}\NormalTok{(}\SpecialCharTok{{-}}\FunctionTok{abs}\NormalTok{(x)) \}}
\end{Highlighting}
\end{Shaded}

The cdf \(F_Y\) is defined as \[
  F_Y(y) = \frac{1}{2}\int_{-\infty}^{y} \exp(-|s|) ds.
\] If \(y \leq 0\), then \[
  F_Y(y) = \frac{1}{2} \int_{-\infty}^{y} \exp(s) ds = \frac{1}{2}\exp(y)
\] and if \(y > 0\), then \[
  F_Y(y) = F_Y(0) + \frac{1}{2} \int_{0}^{y} \exp(-s) ds = \frac{1}{2} + \frac{1}{2}(1 - \exp(-y)),
\] which may be written as the piece-wise function \[
  F_Y(y) =
  \begin{cases}
    \frac{1}{2} \exp(y)                & y \leq 0\\
    1 - \frac{1}{2} \exp(-y) & y > 0.
  \end{cases}
\]

Now, we find the inverse of \(F_Y\). Over \(y \leq 0\), the inverse of
\(F_Y\) is given by \[
  p = \frac{1}{2} \exp(y).
\] Solving for \(y\), we obtain \(y = \log(2p)\). Note that given
\(y \leq 0\), \(p \in (0,0.5]\). Over \(y > 0\), the inverse of \(F_Y\)
is given by \[
  p = 1 - \frac{1}{2} \exp(-y).
\] Solving for \(y\), we obtain \[
  y = -\log(2-2p)
\] where \(p \in (0.5,1)\). We thus have a piecewise function \[
  F^{-1}(p) =
  \begin{cases}
    \log(2p)        & p \in (0,0.5]\\
    -\log(2-2p)     & p \in (0.5,1).
  \end{cases}
\] Thus, to sample from \(Y\), we observe \(u\) from \(U(0,1)\) and, if
\(u \leq 0.5\), we let \(y = \log(2u)\) and otherwise let
\(y = -\log(2-2u)\).

\begin{Shaded}
\begin{Highlighting}[]
\NormalTok{ry }\OtherTok{\textless{}{-}} \ControlFlowTok{function}\NormalTok{(n)}
\NormalTok{\{}
\NormalTok{  ys }\OtherTok{\textless{}{-}} \FunctionTok{numeric}\NormalTok{(n)}
  \ControlFlowTok{for}\NormalTok{ (i }\ControlFlowTok{in} \DecValTok{1}\SpecialCharTok{:}\NormalTok{n)}
\NormalTok{  \{}
\NormalTok{    u }\OtherTok{\textless{}{-}} \FunctionTok{runif}\NormalTok{(}\DecValTok{1}\NormalTok{)}
    \ControlFlowTok{if}\NormalTok{ (u }\SpecialCharTok{\textless{}=} \FloatTok{0.5}\NormalTok{) \{ ys[i] }\OtherTok{\textless{}{-}} \FunctionTok{log}\NormalTok{(}\DecValTok{2}\SpecialCharTok{*}\NormalTok{u) \}}
    \ControlFlowTok{else}\NormalTok{ \{ ys[i] }\OtherTok{\textless{}{-}} \SpecialCharTok{{-}}\FunctionTok{log}\NormalTok{(}\DecValTok{2{-}2}\SpecialCharTok{*}\NormalTok{u)\}}
\NormalTok{  \}}
\NormalTok{  ys}
\NormalTok{\}}
\end{Highlighting}
\end{Shaded}


\end{document}
